A partir de los resultados obtenidos, se concluye que la 
fuerza bruta es inviable más allá de dimensiones de 
aproximadamente $16 \times 16$: 
ya que toma segundos con 256 celdas y su tiempo se multiplica 
por aproximadamente $600$ al duplicar el tamaño. 
La programación dinámica entrega el óptimo hasta aproximadamente$2 \cdot 10^{6}$ 
celdas en aproximadamente $0{,}6$ s, pero usando aproximadamente $230$ MB. 
Greedy es unas $2000$ veces más rápido que DP y casi no usa memoria, pero no 
garantiza el óptimo: su brecha crece linealmente (aproximadamente proporcional a $n+m$)
y falla hasta en el $67\%$ de las instancias con murallas, 
mientras que en 
terrenos libres o con minas casi nunca falla. La mejor elección es, 
por todo lo visto, un compromiso entre garantía de optimalidad, 
memoria y tiempo.